%% =====================================================================
%%  Corrigés — Unit 20 « Happy birthday, Léa! »
%% =====================================================================
\corrigesUnite{20}{Happy birthday, Léa!}

\begin{corrige}{1}
\begin{reponses}[4]
\item \rep{a cake}
\item \rep{a candle}
\item \rep{make a wish}
\item \rep{a present}
\item \rep{a card}
\item \rep{celebrate}
\item \rep{Congratulations!}
\item \rep{Cheers!}
\end{reponses}
\retiens{« un cadeau » se dit \en{a present} ou \en{a gift} : les deux sont corrects.}
\end{corrige}

\begin{corrige}{2}
\begin{reponses}[2]
\item \rep{works}
\item \rep{is making}
\item \rep{bought}
\item \rep{are going to drive}
\item \rep{was}
\item \rep{Did you see}
\item \rep{doesn't like}
\item \rep{is coming}
\end{reponses}
\pourquoi{phrase 1 : un fait habituel (présent simple, \en{-s} de la 3\textsuperscript{e} personne) ;
\en{Look!}, \en{Shh!} : action en cours ; \en{Yesterday}, \en{Six months ago} : prétérit ; \en{Next
summer} : \en{be going to}.}
\end{corrige}

\begin{corrige}{3}
\begin{reponses}[4]
\item \rep{Did}
\item \rep{Can}
\item \rep{Are}
\item \rep{Was}
\item \rep{Does}
\item \rep{Is}
\item \rep{Were}
\item \rep{Do}
\end{reponses}
\pourquoi{regardez le verbe qui suit le sujet : base verbale → \en{do}, \en{does}, \en{did} ou \en{can} ;
verbe en \en{-ing} ou \en{going to} → \en{am}, \en{is}, \en{are} ; adjectif (\en{good}, \en{happy}) →
\en{was}, \en{were}.}
\end{corrige}

\begin{corrige}{4}
\begin{reponses}[3][french]
\item \rep{V}
\item \rep{F}
\item \rep{V}
\item \rep{F}
\item \rep{V}
\item \rep{V}
\end{reponses}
\pourquoi{Léa est arrivée il y a six mois (\en{six months ago}) ; tout le monde aime son travail
(\en{everybody loves your work}).}
\end{corrige}

\begin{corrige}{5}
\begin{reponses}[1]
\item \rep{How old are you?}
\item \rep{Where were you born?}
\item \rep{Where did you go?}
\item \rep{What are you doing?}
\item \rep{What time do you get up?}
\item \rep{Can you swim?}
\end{reponses}
\retiens{la réponse donne le temps de la question : \en{I was born} → \en{were you born?} ; \en{We went}
→ \en{did you go?} ; \en{I'm making} → \en{are you doing?}}
\end{corrige}

\begin{corrige}{6}
\begin{reponses}[1]
\item \rep{I work at Bright Ideas.}
\item \rep{She doesn't like cake.}
\item \rep{Did you go to the party?}
\item \rep{I'm going to travel next year.}
\item \rep{Yesterday I was at home.}
\end{reponses}
\piege{phrase 1 : jamais \en{be} devant une base verbale. \en{I am work} mélange deux temps : \en{I work}
(habitude) ou \en{I'm working} (maintenant).}
\end{corrige}

\begin{corrige}{7}
\begin{reponses}[3][french]
\item \rep{futur}
\item \rep{passé}
\item \rep{présent}
\item \rep{passé}
\item \rep{futur}
\item \rep{présent}
\end{reponses}
\pourquoi{\en{going to} annonce le futur ; \en{bought}, \en{were} sont des formes passées ; \en{live} et
\en{I'm opening} décrivent le présent. Texte de la piste page~\pageref{trans-20.10}.}
\end{corrige}

\begin{corrige}{8}
\begin{reponses}[2]
\item \rep{months}
\item \rep{nervous}
\item \rep{friends}
\item \rep{job}
\item \rep{Hugo}
\end{reponses}
\pourquoi{texte de la piste page~\pageref{trans-20.11}.}
\end{corrige}

\begin{corrige}{9}
\begin{reponses}[3]
\item \rep{Did}
\item \rep{was}
\item \rep{are}
\item \rep{going}
\item \rep{Do}
\item \rep{please}
\item \rep{made}
\item \rep{birthday}
\end{reponses}
\retiens{\en{the best party ever} = « la meilleure fête de tous les temps » ; \en{huge} = « énorme ».}
\end{corrige}

\begin{corrige}{10}
\begin{reponses}[1]
\item \rep{Happy birthday! How old are you?}
\item \rep{Six months ago, I didn't speak English.}
\item \rep{Now I speak English every day.}
\item \rep{Look! They're bringing the cake!}
\item \rep{Next year, I'm going to go to New York.}
\item \rep{Did you like the party? — Yes, I did.}
\item \rep{I'm proud of you!}
\item \rep{Thank you for everything!}
\end{reponses}
\piege{phrase 5 : on peut aussi dire \en{I'm going to New York} : \en{going to go} se simplifie souvent.}
\end{corrige}

\begin{corrige}{11}
Réponse modèle :
\begin{anglais}
\begin{tcolorbox}[otBase,colback=otGreenTint,colframe=otGreenTint,fontupper=\fontsize{10.4}{14}\selectfont]
Six months ago, I didn't speak English. I was very shy. Now I can introduce myself, order in a café
and talk about my weekend. I listen to the audio tracks every day. Next month, I'm going to start
Book 2. And next summer, I'm going to visit London!
\end{tcolorbox}
\end{anglais}
\retiens{\en{introduce myself} = « me présenter ». Bravo pour tout ce chemin parcouru !}
\end{corrige}

\begin{corrige}{12}
Réponses modèles :
\begin{reponses}[1]
\item \rep{My name is Thomas. I'm from Bordeaux, in France.}
\item \rep{I get up at seven and I go to work by bus.}
\item \rep{I visited my parents and I watched a movie.}
\item \rep{I'm listening to Emma and I'm speaking English!}
\item \rep{I'm going to travel to Canada.}
\end{reponses}
\end{corrige}

\begin{corrigetest}
\begin{reponses}[5][french]
\item \rep{b}
\item \rep{a}
\item \rep{c}
\item \rep{a}
\item \rep{b}
\item \rep{c}
\item \rep{a}
\item \rep{b}
\item \rep{c}
\item \rep{a}
\end{reponses}
\pourquoi{questions 5 et 6 : la réponse courte indique l'auxiliaire (\en{did}, \en{does}).}
\end{corrigetest}
